\documentclass[11pt]{article}
\usepackage{ctex}
\usepackage{booktabs}
\usepackage{amsmath}
\usepackage{graphicx}
\title{激光粉末床熔融多孔支架的孔隙率偏差：一个可迁移判据（示例）}
\author{示例作者}
\begin{document}
\maketitle

\begin{abstract}
本文以体心立方与金刚石两类单元结构为对象，在激光功率与扫描速度的二维参数面上界定孔隙率偏差小于百分之四的稳定带，并给出带外的失效模式与适用范围。示例文本，不代表任何真实实验。
\end{abstract}

\section{引言}
多孔支架的强度与渗透性在孔隙率上方向相反，这一取舍在骨修复与轻量化两类场景中都被反复讨论\cite{wang2019}。工程实践因此把孔隙率当作首要设计变量，而非事后调节的工艺副产品：孔隙率每提高一个百分点，渗透性改善的同时承载截面被同步削弱，两个目标共用一个几何自由度。增材制造把孔隙率从配料决定变为结构设计与工艺实现共同决定，可控性却未随自由度同步提升\cite{zhang2020}。

激光粉末床熔融成形点阵时，实测孔隙率与设计值之间的偏差既来自支柱断裂与球化，也来自层间熔合不足\cite{liu2022}。偏差的成因跨越两个尺度：熔池尺度决定单道致密化，结构尺度决定热积累与支撑条件。两个尺度耦合后，偏差不再随单一参数单调漂移，而是在特定功率与速度组合内收敛、在组合之外迅速发散\cite{chen2021}。这意味着"孔隙率可设计"成立的前提，是存在一个可识别的稳定工艺带。

产业中已有若干经验窗口，但多按设备与粉末批次分别标定，缺少可跨结构类型迁移的判据\cite{sun2020}。评价口径同样不统一：有的以致密度代替孔隙率，有的以设计孔隙率为基准计算相对偏差，导致不同来源的数据难以横向比较\cite{park2021}。本文因此把偏差本身作为对象，给出二维参数面上的稳定带边界，并检验该边界是否随单元结构改变。

\section{相关工作}
孔隙率与力学性能的关系研究以 Gibson-Ashby 类模型为主，共同做法是固定拓扑后扫描孔隙率并测力学响应\cite{wang2019}。近期工作把节点形状纳入拟合，发现同一孔隙率下节点加粗可提升刚度约一成，但渗透性同步下降\cite{wu2023}。这一路线的取舍在于模型简洁、可外推到设计阶段，代价是把"孔隙率相同则性能相同"当作前提，未区分孔隙率由设计引入还是由缺陷引入。

工艺参数与缺陷机理研究从熔池行为出发，建立参数与缺陷类型的对应关系并给出致密化窗口\cite{chen2021}。文献\cite{zhang2020}指出支柱几何对缺陷的敏感性随倾角变化，说明窗口边界与结构相关；\cite{rossi2022}用高速摄影证实扫描策略改变热积累分布，同一功率速度组合在分区扫描与连续扫描下致密度可差一点五个百分点。三项工作合起来表明，窗口是参数、结构与扫描策略的联合函数，一维窗口的迁移性因此受限。

渗透性能研究用压降或渗透率实验比较不同单元结构，共同结论是相同孔隙率下不同拓扑的渗透率可相差数十个百分点\cite{kumar2018}。综述性工作梳理了窗口概念，但仍以定性描述为主\cite{liu2022}；\cite{zheng2023}尝试用孔道迂曲度统一不同拓扑的渗透数据，在两类结构上拟合良好，但样本仅覆盖致密区间。选型结论的适用边界恰好落在带外区域，而该区域尚无系统数据。

\section{方法}
参数扫描覆盖激光功率一百二十至二百瓦、扫描速度六百至一千二百毫米每秒，层厚三十微米，材料为 Ti-6Al-4V，粉末粒径十五至四十五微米。两类单元结构的支柱直径保持一致，以隔离拓扑差异；扫描策略统一为条带分区，条带宽度五毫米，层间旋转六十七度以抑制热积累的方向性。

孔隙率由阿基米德法实测，每个参数点重复三次取均值，标准差超过百分之零点五的点重测。偏差定义为实测值与设计值之差除以设计值。稳定带判据写成
\begin{equation}
\delta(P,v) = \frac{\phi_{\mathrm{exp}}(P,v) - \phi_{\mathrm{des}}}{\phi_{\mathrm{des}}}, \qquad |\delta| < 0.04,
\end{equation}
其中 $P$ 为激光功率，$v$ 为扫描速度。判据的作用是把"可造"从定性判断变为在二维面上可计算的条件：给定结构与参数，$\delta$ 是否落在阈值内可以直接判定。带外样本按球化与塌陷两类失效模式分别归类，归类依据为成形面宏观形貌与截面致密度分布。

为使判据可迁移，偏差的口径在全部参数点复用同一测量流程与同一基准（设计孔隙率），不使用致密度代替孔隙率，也不按设备单独标定阈值。该选择使不同结构、不同参数点的数值可横向比较，代价是对测量误差更敏感，因此每个参数点都保留三次重复的原始值。

\section{结果与讨论}
表~\ref{tab:main} 给出两类结构在代表性参数点上的偏差与断裂率。稳定带内偏差普遍低于百分之四，孔隙率超过百分之七十二后体心立方结构偏差升至百分之五至九，金刚石结构的支柱断裂率由百分之四升至百分之十二，但渗透率高出约百分之四十。该优势仅在稳定带内可复现。

\begin{table}[htbp]
\centering
\caption{两类单元结构的孔隙率偏差与支柱断裂率（示例数据）}
\label{tab:main}
\begin{tabular}{lccc}
\toprule
结构 & 孔隙率(\%) & 偏差(\%) & 断裂率(\%) \\
\midrule
BCC & 65 & 2.1 & 4 \\
BCC & 74 & 6.8 & 5 \\
Diamond & 66 & 2.9 & 9 \\
Diamond & 75 & 8.4 & 12 \\
\bottomrule
\end{tabular}
\end{table}

带外失效并非随机分布，而是沿参数面呈现两条带状区域：低功率高速侧以未熔合为主，高功率低速侧以球化与塌陷为主。两条带的边界与单元结构的支柱倾角相关，金刚石结构的低功率侧边界更靠内，这解释了同等参数下其断裂率更高的现象。稳定带的宽度在两类结构间相差约百分之十五，说明边界确实随拓扑改变，判据不能跨结构直接套用。

\section{结论}
孔隙率偏差在二维参数面上存在可识别的稳定带，其边界随单元结构改变：体心立方结构在孔隙率超过百分之七十二后偏差明显放大，金刚石结构的支柱断裂率在同等参数下更高。稳定带宽度在两类结构间相差约百分之十五，带外失效沿参数面呈两条带状分布，低功率高速侧以未熔合为主、高功率低速侧以球化与塌陷为主。把渗透与力学的比较限制在带内，两类结论才可复现；带外的比较属于失效模式描述，不构成适用范围。就工程使用而言，判据给出的是"该结构在该参数下是否可造"的二值回答，而不是对性能的预测，因此不能替代性能测试。本文结论不外推到其他合金体系，也不适用于层厚显著偏离三十微米或扫描策略改为连续扫描的情形；后者的热积累分布不同，稳定带边界需重新标定。

\bibliographystyle{plain}
\bibliography{selftest-refs}
\end{document}
